\runninghead{The Appointed Texts}
\section{The Appointed Texts}
\label{sec:appointed-texts}

The Latin below is that of the 1962 typical edition, \work{Missale Romanum}
(Typis Polyglottis Vaticanis), printed pp.~1--2, nos.~1--10, read on the page
images of the CMAA facsimile, where the Mass stands first in the book under
the heading \latin{Dominica prima Adventus}, I classis, \latin{Statio ad S.
Mariam maiorem}. Every word of the ten texts stands in the Ratisbon printing
of the \work{Missale Romanum} (Pustet, 1862), pp.~1--3. The two printings
differ in spelling, pointing and capitals, in the form of the citations and in
the conclusion cues, and in no word that is said.

The English of the three orations is that of the Lasance \work{New Roman
Missal} in the Benziger revision of 1945, which prints this Mass in Latin and
English under the heading ``First Sunday of Advent.'' The English of the
lessons and chants is the Douay--Rheims Bible in Bishop Challoner's revision,
which translates the Clementine Vulgate; each place where the Missal's words
depart from the Vulgate is named beneath the text. For the Communion the
Lasance line is given, since it renders the Missal's shorter form. None of
this English is an approved liturgical translation of the 1962 Missal. Psalms
are cited in the Missal's numbering, with the Hebrew number in parentheses.

\subsection{Introit}
\label{proper-introit}

\properlocus{\latin{Antiphona ad Introitum}, \latin{Ps.\ 24, 1--3} (Heb.\ 25);
verse \latin{Ps.\ 24, 4}. Missal p.~1, no.~1.}

\begin{latinproper}
Ad te levávi ánimam meam: Deus meus, in te confído, non erubéscam: neque
irrídeant me inimíci mei: étenim univérsi, qui te exspéctant, non
confundéntur. \textnormal{Ps.\ ibid., 4} Vias tuas, Dómine, demónstra mihi:
et sémitas tuas édoce me. ℣.~Glória Patri, et Fílio, et Spirítui Sancto. Sicut
erat in princípio, et nunc, et semper, et in sǽcula sæculórum. Amen.
\end{latinproper}

\begin{englishwitness}{Douay--Rheims (Challoner), Ps 24:1--3, the antiphon}
To thee, O Lord, have I lifted up my soul. In thee, O my God, I put my trust;
let me not be ashamed. Neither let my enemies laugh at me: for none of them
that wait on thee shall be confounded.
\end{englishwitness}

\begin{englishwitness}{Douay--Rheims (Challoner), Ps 24:4b, the verse}
Shew, O Lord, thy ways to me, and teach me thy paths.
\end{englishwitness}

\noindent The antiphon opens \latin{Ad te levávi} without the vocative of the
Vulgate's \latin{Ad te, Domine, levavi}, so the Douay's ``O Lord'' in the first
line has no word in the chant. Where the Vulgate reads \latin{qui sustinent te},
which the Douay renders ``that wait on thee,'' the chant reads \latin{qui te
exspéctant}, ``who wait for thee.'' The verse agrees with the Vulgate. After
the doxology the antiphon is repeated.

\subsection{Collect}
\label{proper-collect}

\properlocus{\latin{Oratio}. Missal p.~1, no.~2. The Gloria is not said in
Masses of the season from this Sunday to the Vigil of the Nativity.}

\begin{latinproper}
Excita, quǽsumus, Dómine, poténtiam tuam, et veni: ut ab imminéntibus
peccatórum nostrórum perículis, te mereámur protegénte éripi, te liberánte
salvári: Qui vivis et regnas cum Deo Patre in unitáte Spíritus Sancti, Deus:
per ómnia sǽcula sæculórum. R.~Amen.
\end{latinproper}

\begin{englishwitness}{Lasance, The New Roman Missal (1945), First Sunday of Advent, Collect}
Bestir, O Lord, Thy might, we pray Thee and come; That, defended by Thee, we
may deserve rescue from approaching dangers brought on by our sins, and being
set free by Thee, obtain our salvation. Who livest.
\end{englishwitness}

\noindent The conclusion \latin{Qui vivis et regnas cum Deo Patre} shows that
the prayer is addressed to the Son.

\subsection{Epistle}
\label{proper-epistle}

\properlocus{\latin{Léctio Epístolæ beáti Pauli Apóstoli ad Romános},
\latin{Rom.\ 13, 11--14}. Missal pp.~1--2, no.~3.}

\begin{latinproper}
Fratres: Sciéntes, quia hora est iam nos de somno súrgere. Nunc enim própior
est nostra salus, quam cum credídimus. Nox præcéssit, dies autem
appropinquávit. Abiciámus ergo ópera tenebrárum, et induámur arma lucis.
Sicut in die honéste ambulémus: non in comessatiónibus, et ebrietátibus, non
in cubílibus, et impudicítiis, non in contentióne, et æmulatióne: sed
induímini Dóminum Iesum Christum.
\end{latinproper}

\begin{englishwitness}{Douay--Rheims (Challoner), Rom 13:11--14a}
And that, knowing the season, that it is now the hour for us to rise from
sleep. For now our salvation is nearer than when we believed. The night is
passed And the day is at hand. Let us, therefore cast off the works of
darkness and put on the armour of light. Let us walk honestly, as in the day:
not in rioting and drunkenness, not in chambering and impurities, not in
contention and envy. But put ye on the Lord Jesus Christ.
\end{englishwitness}

\noindent The Missal opens the lesson \latin{Fratres: Sciéntes}, which Lasance
renders ``Brethren, knowing the time,'' where the Vulgate has \latin{Et hoc
scientes tempus}, ``And that, knowing the season.'' The lesson ends at the
first half of verse~14. The rest of the verse, ``and make not provision for
the flesh in its concupiscences,'' is not read.

\subsection{Gradual}
\label{proper-gradual}

\properlocus{\latin{Graduale}, \latin{Ps.\ 24, 3 et 4}. Missal p.~2, no.~4.}

\begin{latinproper}
Univérsi, qui te exspéctant, non confundéntur, Dómine. ℣.~Vias tuas, Dómine,
notas fac mihi: et sémitas tuas édoce me.
\end{latinproper}

\begin{englishwitness}{Douay--Rheims (Challoner), Ps 24:3b--4b}
For none of them that wait on thee shall be confounded. \dots\ Shew, O Lord,
thy ways to me, and teach me thy paths.
\end{englishwitness}

\noindent The respond ends with a \latin{Dómine} that the psalm verse does not
have and that neither the Douay nor Lasance renders, and it reads \latin{qui
te exspéctant} as the Introit does. The verse asks God to ``make known'' his
ways, \latin{notas fac mihi}, where the Introit verse and the Vulgate have
\latin{demónstra mihi}, ``shew to me''; the Douay above translates the
Vulgate's wording. On the ferias of Advent, when this Mass is repeated, the
Gradual is said and the Alleluia with its verse is omitted.

\subsection{Alleluia}
\label{proper-alleluia}

\properlocus{\latin{Ps.\ 84, 8} (Heb.\ 85). Missal p.~2, no.~5.}

\begin{latinproper}
Allelúia, allelúia. ℣.~Osténde nobis, Dómine, misericórdiam tuam: et
salutáre tuum da nobis. Allelúia.
\end{latinproper}

\begin{englishwitness}{Douay--Rheims (Challoner), Ps 84:8}
Shew us, O Lord, thy mercy; and grant us thy salvation.
\end{englishwitness}

\noindent The verse agrees with the Vulgate.

\subsection{Gospel}
\label{proper-gospel}

\properlocus{\latin{Sequéntia sancti Evangélii secúndum Lucam},
\latin{Luc.\ 21, 25--33}. Missal p.~2, no.~6. The Credo is said.}

\begin{latinproper}
In illo témpore: Dixit Iesus discípulis suis: Erunt signa in sole, et luna,
et stellis, et in terris pressúra géntium præ confusióne sónitus maris, et
flúctuum: arescéntibus homínibus præ timóre et exspectatióne, quæ
supervénient univérso orbi: nam virtútes cælórum movebúntur. Et tunc vidébunt
Fílium hóminis veniéntem in nube cum potestáte magna, et maiestáte. His autem
fíeri incipiéntibus, respícite, et leváte cápita vestra: quóniam appropínquat
redémptio vestra. Et dixit illis similitúdinem: Vidéte ficúlneam, et omnes
árbores: cum prodúcunt iam ex se fructum, scitis quóniam prope est æstas. Ita
et vos cum vidéritis hæc fíeri, scitóte quóniam prope est regnum Dei. Amen
dico vobis, quia non præteríbit generátio hæc, donec ómnia fiant. Cælum et
terra transíbunt: verba autem mea non transíbunt.
\end{latinproper}

\begin{englishwitness}{Douay--Rheims (Challoner), Lk 21:25--33}
And there shall be signs in the sun and in the moon and in the stars; and upon
the earth distress of nations, by reason of the confusion of the roaring of
the sea, and of the waves: Men withering away for fear and expectation of
what shall come upon the whole world. For the powers of heaven shall be
moved. And then they shall see the Son of man coming in a cloud, with great
power and majesty. But when these things begin to come to pass, look up and
lift up your heads, because your redemption is at hand. And he spoke to them
a similitude. See the fig tree and all the trees: When they now shoot forth
their fruit, you know that summer is nigh; So you also, when you shall see
these things come to pass, know that the kingdom of God is at hand. Amen, I
say to you, this generation shall not pass away till all things be fulfilled.
Heaven and earth shall pass away: but my words shall not pass away.
\end{englishwitness}

\noindent The Missal opens the lesson \latin{In illo témpore: Dixit Iesus
discípulis suis: Erunt signa}, which Lasance renders ``At that time, Jesus
said to His disciples: There shall be signs,'' where the Vulgate has \latin{Et
erunt signa}, ``And there shall be signs.'' In St Luke that \latin{Et} follows
the siege and fall of Jerusalem (Lk 21:20--24). The rest of the lesson agrees
with the Vulgate.

\subsection{Offertory}
\label{proper-offertory}

\properlocus{\latin{Antiphona ad Offertorium}, \latin{Ps.\ 24, 1--3}. Missal
p.~2, no.~7.}

\begin{latinproper}
Ad te levávi ánimam meam: Deus meus, in te confído, non erubéscam: neque
irrídeant me inimíci mei: étenim univérsi, qui te exspéctant, non
confundéntur.
\end{latinproper}

\begin{englishwitness}{Douay--Rheims (Challoner), Ps 24:1--3}
To thee, O Lord, have I lifted up my soul. In thee, O my God, I put my trust;
let me not be ashamed. Neither let my enemies laugh at me: for none of them
that wait on thee shall be confounded.
\end{englishwitness}

\noindent The words are those of the Introit's antiphon, with the same two
departures from the Vulgate.

\subsection{Secret}
\label{proper-secret}

\properlocus{\latin{Secreta}. Missal p.~2, no.~8. The Preface of the Most
Holy Trinity follows.}

\begin{latinproper}
Hæc sacra nos, Dómine, poténti virtúte mundátos, ad suum fáciant purióres
veníre princípium. Per Dóminum nostrum Iesum Christum, Fílium tuum: Qui tecum
vivit et regnat in unitáte Spíritus Sancti, Deus: per ómnia sǽcula
sæculórum. R.~Amen.
\end{latinproper}

\begin{englishwitness}{Lasance, The New Roman Missal (1945), First Sunday of Advent, Secret}
Cleansing us by their mighty power, may these Holy Mysteries, O Lord, make
us come more pure before Thee who art their author. Through our Lord.
\end{englishwitness}

\noindent The Latin asks that the holy things make us come purer \latin{ad suum
princípium}, ``to their own beginning''; Lasance renders the phrase ``before
Thee who art their author.''

\subsection{Communion}
\label{proper-communion}

\properlocus{\latin{Antiphona ad Communionem}, \latin{Ps.\ 84, 13} (Heb.\ 85).
Missal p.~2, no.~9.}

\begin{latinproper}
Dóminus dabit benignitátem: et terra nostra dabit fructum suum.
\end{latinproper}

\begin{englishwitness}{Lasance, The New Roman Missal (1945), First Sunday of Advent, Communion}
The Lord will give goodness: and our earth shall yield her fruit.
\end{englishwitness}

\begin{englishwitness}{Douay--Rheims (Challoner), Ps 84:13}
For the Lord will give goodness: and our earth shall yield her fruit.
\end{englishwitness}

\noindent The Vulgate's verse begins \latin{Etenim}, which the Douay renders
``For''; the chant begins with \latin{Dóminus}, as Lasance's line does.

\subsection{Postcommunion}
\label{proper-postcommunion}

\properlocus{\latin{Postcommunio}. Missal p.~2, no.~10.}

\begin{latinproper}
Suscipiámus, Dómine, misericórdiam tuam in médio templi tui: ut reparatiónis
nostræ ventúra solémnia cóngruis honóribus præcedámus. Per Dóminum nostrum.
\end{latinproper}

\begin{englishwitness}{Lasance, The New Roman Missal (1945), First Sunday of Advent, Postcommunion}
May we receive Thy mercy, O Lord, in the midst of Thy temple, that with due
reverence we may prepare for the coming festival of our redemption. Through
our Lord.
\end{englishwitness}

\noindent The opening words are those of Ps 47:10 (48:10), \latin{Suscepimus,
Deus, misericordiam tuam in medio templi tui}, ``We have received thy mercy,
O God, in the midst of thy temple,'' turned from the past tense to a request.
The Latin's \latin{præcedámus} is ``go before''; its \latin{reparatiónis
nostræ} is ``of our restoration,'' which Lasance renders ``of our redemption.''
